\newcommand{\HRule}[1]{\rule{\linewidth}{#1}}

\declaretheoremstyle[name=Theorem,]{thmsty}
\declaretheorem[style=thmsty,numberwithin=section]{theorem}
\tcolorboxenvironment{theorem}{colback=LightGray}

\declaretheoremstyle[name=Proposition,]{prosty}
\declaretheorem[style=prosty,numberlike=theorem]{proposition}
\tcolorboxenvironment{proposition}{colback=LightOrange}

\declaretheoremstyle[name=Principle,]{prcpsty}
\declaretheorem[style=prcpsty,numberlike=theorem]{principle}
\tcolorboxenvironment{principle}{colback=LightGreen}

\newcommand{\mysection}[2]{%
    \setcounter{section}{#2}
    \section*{#1}%
    \addcontentsline{toc}{section}{#1}%
    \markright{#1}%
}

\newcommand{\mysubsection}[2]{%
    \setcounter{subsection}{#2}
    \subsection{#1}%
    \markright{#1}%
}

\newcommand{\contradiction}{\ensuremath{{\Rightarrow\mspace{-2mu}\Leftarrow}}}
\newcommand{\myquestion}[1]{\noindent \textbf{P.#1)}}

\newcommand{\ip}{{\color{red}(\textit{\textbf{IN PROGRESS}})}\;}
\newcommand{\fp}{{\color{red}(\texttt{1st pass})}\;}
\newcommand{\sndp}{{\color{blue}(\texttt{2nd pass})}\;}
\newcommand{\tp}{{\color{green}(\texttt{3rd pass})}\;}

% General Commands
\newcommand{\inv}[1]{#1^{-1}}
\newcommand{\norm}[1]{\left\lVert#1\right\rVert}
\newcommand{\abs}[1]{\left|#1\right|}
\newcommand{\paren}[1]{\left(#1\right)}
\newcommand{\brak}[1]{\left\langle#1\right\rangle}
\newcommand{\incomplete}[1]{\noindent \textcolor{red}{#1}}
\newcommand{\textbt}[1]{\textit{\textbf{#1}}}
\newcommand{\binrep}[1]{{}_{\llcorner\mkern-3mu}{#1}_{\mkern-3mu\lrcorner}}

% Distributions
\newcommand{\Unif}{\text{Unif}}
\newcommand{\Poisson}{\text{Poisson}}
\newcommand{\Gam}{\text{Gamma}}
\newcommand{\Normal}[2]{\text{Norm}(#1, #2)}

% General Math
\newcommand{\deriv}[2]{\frac{d #1}{d #2}}
\newcommand{\pderiv}[2]{\frac{\partial #1}{\partial #2}}

% Turing Machine
\newcommand{\start}{\triangleright}
\newcommand{\blank}{\square}

% Complexity Classes
\renewcommand{\P}{\mathcal{P}}
\newcommand{\NP}{\mathcal{NP}}
\newcommand{\redto}{\leq_p}
\newcommand{\NPhard}{\mathcal{NP}\text{-hard}}
\newcommand{\NPcomplete}{\mathcal{NP}\text{-complete}}
\newcommand{\dtime}[1]{\mathbf{DTIME}(#1)}
\newcommand{\ntime}[1]{\mathbf{NTIME}(#1)}

\newcommand{\bigoh}[1]{\mathcal{O}(#1)}
\newcommand{\littleoh}[1]{\mathcal{o}(#1)}
\newcommand{\bigtheta}[1]{\Theta(#1)}
\newcommand{\bigomega}[1]{\Omega(#1)}
\newcommand{\littleomega}[1]{\omega(#1)}

% Environments
\newmdenv[
  linecolor=black, % Frame color
  innerleftmargin=8pt, % Padding inside the frame
  innerrightmargin=8pt,
  innertopmargin=10pt,
  innerbottommargin=10pt,
  linewidth=.5pt, % Frame line
  frametitlerule=true, %set title true
  frametitlerulewidth=.25pt,%width
]{thoughtbox}

\newenvironment{pencil}[1]
  {
    \par\noindent #1

    \medskip

    \textbf{Proof:}
  }
  {
    \bigskip
  }

\newenvironment{question}[2]
  {
    \par\noindent\myquestion{#1} #2

    \medskip

    \textbf{SOLUTION:}
  }
  {
    \bigskip
  }

\newenvironment{thoughts}[1]
  {
    \begin{thoughtbox}[frametitle={Thoughts:~#1}]
  }
  {\end{thoughtbox}
    \medskip}

\newenvironment{notes}
  {
    \par\noindent\textbf{\uline{Notes}}
    \begin{itemize}[label={\textbullet}, topsep = -.5pt, leftmargin=*]
  }
  {\end{itemize}
    \bigskip}

\newenvironment{enumalph}
  {\begin{enumerate}[label=\alph*), topsep=-0.5pt, itemsep=.5pt]}
  {\end{enumerate}}

\renewenvironment{bullet}
  {\begin{itemize}[topsep=-0.5pt, itemsep=.5pt]}
  {\end{itemize}}

